\section{Conclusion}
In this work, we propose \textbf{Context-as-Memory}, highlighting that using historical generated frames as memory is key to achieving scene-consistent long video generation. Our method design is simple yet effective, directly saving context frames as memory and inputting the context together with the predicted frame as conditions. Furthermore, to avoid high computational overhead caused by lengthy context, we propose \textbf{Memory Retrieval} to dynamically select truly valuable context based on the predicted video frames. 
\paragraph{Limitations and Future Work.}
Although our method has made significant progress in achieving memory capability for long video generation, several limitations remain: (1) Our method is limited to static scenes, while memory retrieval for dynamic scenes poses greater challenges; (2) In complex scenarios, particularly those with multiple occlusions (e.g., interconnected indoor rooms), FOV overlap may struggle to effectively identify truly relevant context frames; (3) The inherent error accumulation problem in long video generation persists, which currently can only be addressed through larger datasets, more extensive training, and more powerful base models. In the future, we will continue to develop memory capabilities for open-domain long video generation on larger-scale base models, supporting more complex trajectories, broader scene ranges, and longer generation sequences.